\section{Study Design}

\subsection{Overall Design}
A prospective, open-label, single-arm, first-in-human study combining a Phase 1
3+3 dose-finding evaluation of perioperative daraxonrasib with an early
feasibility evaluation of the on-premises LLM-directed eight-arm robotic Whipple
(significant-risk device, Class II collaborative).

\draftinstr{State the hypothesis, the combined IND/IDE phase designation, the
single-arm rationale, n up to 18, and the staged-autonomy frame; render
\texttt{mermaid/fig-18-staged-autonomy.md}. Source the Class II / continuous
oversight rules from
\texttt{inputs/21cfr312\_adapt/05\_clinical\_holds\_appendices\_closing.tex}.}

\subsection{Scientific Rationale for Study Design}
\draftinstr{Justify single-arm first-in-human (vs randomized) by the
device-feasibility precedent and the vulnerable population; cite
\texttt{research/ChatGPT-5.5-Thinking-Extended-19Jun26.md}.}

\subsection{Justification for Dose}
\draftinstr{Justify daraxonrasib starting dose and 3+3 escalation in exact doses
(not percentages) anchored to RASolute 302; render
\texttt{mermaid/fig-10-dose-escalation.md}. Justify the device ``dose'' analogue:
USL >= 7.0 surgical readiness, Phase 0 >= 1000 simulated procedures across >= 2
frameworks, sim-to-real < 2mm / < 0.5N, from
\texttt{inputs/21cfr312\_adapt/02\_ind\_content\_phases.tex}.}

\subsection{End of Study Definition}
\draftinstr{Adapt the NIH End-of-Study language: a participant completes the
study at the last visit / procedure in the SoA; the end of study is the last
participant's 24-month OS assessment. Source from
\texttt{nih-protocol/03\_study\_design\_and\_study\_population.md}.}
